\input{../tex-inputs/latex-leseansicht-vorspann}

\section[Richard Beer-Hofmann an Arthur Schnitzler, 14. 9. 1896]{L00589 Richard Beer-Hofmann an Arthur Schnitzler, 14. 9. 1896}
\nopagebreak\mylabel{L00589v}
\rehead{ }\normalsize\beginnumbering\briefempfaengerindex{Schnitzler, Arthur@\textsc{Schnitzler, Arthur}!zzzBeer-Hofmann, Richard@\emph{von Richard Beer-Hofmann}!1896-09-141@{14. 9. 1896}|(be}
\toendnotes[C]{\smallbreak\pagebreak[2]}
\correspDesc{Versand  durch Richard Beer-Hofmann am 14. 9. 1896 in Kaiser-Franz-Ring 54
\newline{}Übermittlung  am 14. 9. 1896 in Baden bei Wien
\newline{}Zustellung  am 14. 9. 1896 in 9/3 Post- und Telegraphenamt 72
\newline{}Erhalt  durch Arthur Schnitzler am 14. 9. 1896 in Frankgasse 1}\toendnotes[C]{\smallbreak}
\Standort{CUL, Schnitzler, B 8.}
\physDesc{Postkarte, 311 Zeichen
\newline{}Handschrift: Bleistift, lateinische Kurrent
\newline{}Versand: 1) Stempel: »\nobreak{}\oindex{Baden bei Wien@\textbf{Baden bei Wien}|pwk}Baden 1 N. Ö., 14. 9. 96, 2–3 N\nobreak{}«.   2) Stempel: »\nobreak{}\oindex{Wien@\textbf{Wien}!IX., Alsergrund@\textbf{IX., Alsergrund}!9/3 Post- und Telegraphenamt 72@\textbf{9/3 Post- und Telegraphenamt 72}|pwk}Wien 9/3 72, 14. 9. 96, 7.N, Bestellt\nobreak{}«. 
\newline{}Ordnung: mit Bleistift von unbekannter Hand nummeriert:
                                    »84« }
\buchAbdrucke{\weitereDrucke{Arthur Schnitzler, Richard Beer-Hofmann: \emph{Briefwechsel 1891–1931}. Herausgegeben von Konstanze Fliedl. Wien, Zürich: \emph{Europaverlag} 1992, S. 96.} }\toendnotes[C]{\smallbreak}\pstart{}{\pb}Herrn\pend{}\pstart{}D\textsuperscript{r} Arthur Schnitzler\pend{}\pstart{}Wien\oindex{Wien@\textbf{Wien}|pw}\pend{}\pstart{}IX Frankgasse 1\oindex{Wien@\textbf{Wien}!IX., Alsergrund@\textbf{IX., Alsergrund}!Frankgasse 1@\textbf{Frankgasse 1}|pw}\pend{}{\bigskip}\vspace{1em}
\pstart
           \noindent{}{\pb}Lieber Arthur! Ich war \label{K_L00589-1v}\edtext{gestern}{\lemma{\textnormal{\emph{gestern}}}\Cendnote{\textnormal{Vgl. A. S.: \emph{Wiener Schnitzler}, 13. 9. 1896.}}}\label{K_L00589-1} den ganzen Nachmittag
               bis ½ 10 Nachts{ }Franzensstraße 54\oindex{Kaiser-Franz-Ring 54@\textbf{Kaiser-Franz-Ring 54}|pw} – allerdings hinter einer
               versperrten Doppelthüre. Daß wir Sie nicht Klopfen gehört haben ist räthselhaft. Ich
               dürfte am 24. in Wien\oindex{Wien@\textbf{Wien}|pw} sein. Sehe ich
               Sie noch vorher?\pend
           
\pstart
           Herzlichst{\\[\baselineskip]}Ihr{\\[\baselineskip]}\spacefill\mbox{Richard}\pend
           \leftskip=0em{}
\pstart
           14/IX 96\pend
           \selectlanguage{ngerman}\endnumbering\briefempfaengerindex{Schnitzler, Arthur@\textsc{Schnitzler, Arthur}!zzzBeer-Hofmann, Richard@\emph{von Richard Beer-Hofmann}!1896-09-141@{14. 9. 1896}|)be}\mylabel{L00589h}  \newcommand{\dateiname}{L00589}\newcommand{\titel}{Richard Beer-Hofmann an Arthur Schnitzler, 14. 9. 1896}\newcommand{\editorInnen}{Martin Anton Müller und Gerd-Hermann Susen}\input{../tex-inputs/latex-leseansicht-abspann}
